\documentclass[10pt,letterpaper]{article}
\usepackage[T1]{fontenc}
\usepackage[utf8]{inputenc}
\usepackage{lmodern}
\usepackage[margin=0.76in]{geometry}
\usepackage[table]{xcolor}
\usepackage{array,tabularx,longtable,booktabs}
\usepackage{enumitem}
\usepackage{amssymb,amsmath}
\usepackage{hyperref}
\usepackage{fancyhdr,lastpage}
\usepackage{microtype}
\usepackage{parskip}
\definecolor{cdblue}{HTML}{203957}
\definecolor{cdgray}{HTML}{586677}
\hypersetup{colorlinks=true,linkcolor=cdblue,urlcolor=cdblue,pdftitle={Prerequisite-Based Continuous Delivery Curriculum}}
\pagestyle{fancy}
\fancyhf{}
\lhead{\small\color{cdgray}Continuous Delivery Engineering}
\rhead{\small\color{cdgray}Curriculum}
\cfoot{\thepage\ / \pageref{LastPage}}
\setlength{\headheight}{14pt}
\setlength{\emergencystretch}{2em}
\setlist[itemize]{leftmargin=*,itemsep=3pt,topsep=4pt}
\setlist[enumerate]{leftmargin=*,itemsep=3pt,topsep=4pt}
\begin{document}
\begin{center}
{\LARGE\bfseries\color{cdblue}Prerequisite-Based Continuous Delivery Curriculum}\\[7pt]
{\large Six stages across sixteen weeks}\\[6pt]
{\small Reference edition: 2026-10-30 $\bullet$ Original educational and operational material}
\end{center}
\noindent\color{cdblue}\rule{\linewidth}{1.2pt}\color{black}
\vspace{0.4em}
\section{Entry assessment}
\begin{itemize}
\item Write a small executable test and interpret a failing assertion.
\item Explain an immutable Git commit versus a mutable image tag.
\item Show ability to read a Terraform plan and AWS IAM trust statement.
\item Describe risk of executing untrusted pull request code with production credentials.
\end{itemize}

\section{1. Foundations and source integrity}
Duration: Weeks 1--2. Prerequisite: Git, shell, Python fundamentals. Primary reading: Chapters 1--3.

\subsection{Mastery evidence}
\begin{itemize}
\item Versioned app, branch protection proposal, CI baseline
\item Peer-reviewed design rationale including an explicit alternative and trade-off.
\item Repeatable demonstration with a documented negative or recovery case.
\end{itemize}

Exit requirement: score at least 16/20 on the associated artifact, and close any safety-critical failure before advancing.


\section{2. Quality and executable acceptance}
Duration: Weeks 3--4. Prerequisite: Phase 1 pass. Primary reading: Chapters 4, 7, 8.

\subsection{Mastery evidence}
\begin{itemize}
\item Unit/contract/acceptance portfolio with negative cases
\item Peer-reviewed design rationale including an explicit alternative and trade-off.
\item Repeatable demonstration with a documented negative or recovery case.
\end{itemize}

Exit requirement: score at least 16/20 on the associated artifact, and close any safety-critical failure before advancing.


\section{3. Artifact pipeline and traceability}
Duration: Weeks 5--6. Prerequisite: Phase 2 pass. Primary reading: Chapters 5--6, 13.

\subsection{Mastery evidence}
\begin{itemize}
\item Single-build artifact, digest and promotion manifest
\item Peer-reviewed design rationale including an explicit alternative and trade-off.
\item Repeatable demonstration with a documented negative or recovery case.
\end{itemize}

Exit requirement: score at least 16/20 on the associated artifact, and close any safety-critical failure before advancing.


\section{4. Terraform and AWS deployment}
Duration: Weeks 7--9. Prerequisite: Phase 3 pass, AWS sandbox. Primary reading: Chapters 10--11, 14.

\subsection{Mastery evidence}
\begin{itemize}
\item OIDC trust, isolated IaC, rehearsed deploy and rollback
\item Peer-reviewed design rationale including an explicit alternative and trade-off.
\item Repeatable demonstration with a documented negative or recovery case.
\end{itemize}

Exit requirement: score at least 16/20 on the associated artifact, and close any safety-critical failure before advancing.


\section{5. Data, security and progressive delivery}
Duration: Weeks 10--12. Prerequisite: Phase 4 pass. Primary reading: Chapters 9, 10, 12.

\subsection{Mastery evidence}
\begin{itemize}
\item Schema compatibility matrix, policy gates, canary test
\item Peer-reviewed design rationale including an explicit alternative and trade-off.
\item Repeatable demonstration with a documented negative or recovery case.
\end{itemize}

Exit requirement: score at least 16/20 on the associated artifact, and close any safety-critical failure before advancing.


\section{6. Governance and capstone}
Duration: Weeks 13--16. Prerequisite: All prior phases. Primary reading: Chapters 15 and cross-chapter review.

\subsection{Mastery evidence}
\begin{itemize}
\item Capstone evidence dossier, metrics baseline and scorecard
\item Peer-reviewed design rationale including an explicit alternative and trade-off.
\item Repeatable demonstration with a documented negative or recovery case.
\end{itemize}

Exit requirement: score at least 16/20 on the associated artifact, and close any safety-critical failure before advancing.


\section{Dependency DAG}
\begin{verbatim}
Git + programming + testing
    |
CI discipline -> test strategy -> immutable package -> promotion pipeline
                                              |
                                   AWS OIDC + Terraform
                                              |
                              progressive release + data safety
                                              |
                                 evidence + maturity review
\end{verbatim}

\section{Assessment policy}
Practice runs are non-production. Never deploy cloud resources without verifying cost, account, region, IAM role, Terraform state isolation, permissions, and cleanup. An advanced learner may demonstrate mastery by supplying equivalent evidence, not by bypassing safety gates.



\end{document}
